\section{Related Work}
\paragraph{Behavioral and temporal descriptions.}
\citet{gandhi2025habits} study verification, backtracking, subgoals, and backward chaining in self-improving reasoners. \citet{kim2026society} connect questions, perspective shifts, conflict, and reconciliation to learning and intervention. We use both vocabularies as overlapping annotations, without assuming independent latent systems. LOT studies reasoning patterns and test-time style alignment \citep{chen2025lot}; ReasonOps induces recurring operations across models and benchmarks \citep{lee2026reasonops}; ThinkARM analyzes functional episodes under efficiency interventions \citep{li2025thinkarm}. Our focus is the connection between within-model domain differences, matched task demands, and the consequences of changing profiles across objective and open tasks.

\paragraph{Task organization and cognitive control.}
\citet{minsky1986society} emphasizes the organization of simpler processes. Human frontoparietal connectivity changes with task demands \citep{cole2013hubs}, and primate population dynamics support context-dependent computation in shared circuitry \citep{mante2013context}. These findings motivate examining organization as well as operation frequency. Expected-value-of-control and resource-rational accounts relate cognitive allocation to benefits and costs \citep{shenhav2013control,lieder2020resource}; \citet{desabbata2024metareasoning} apply metareasoning to LLM computation. Our budget comparisons connect to this perspective without establishing optimal allocation or neural homology.

\paragraph{Steering, transfer, and reward learning.}
SEAL evaluates reasoning steering, efficiency, and transfer \citep{chen2025seal}; the Halt Vector studies a stopping-related intervention and its incorporation into model weights \citep{jayabahu2026halt}. \citet{sinii2025vectors} train residual steering vectors with RL and examine their mechanisms. We hold the profile-to-activation construction fixed across source domains, distinguish domain means from archetypes, and compare natural with forced answering. The RL comparison concerns convergence and differences in measured behavior; bias-only optimization is distinct from full-model post-training.

\paragraph{Creative cognition.}
Human creative idea production involves dynamic cooperation between default and executive networks \citep{beaty2015creative}. \citet{ismayilzada2026creative} examine LLM--brain representational alignment during creative thinking and differences associated with post-training. We study the behavioral consequences of applying interventions across objective and creative tasks, including answer availability, originality, and feasibility. These are functional comparisons; we do not infer corresponding brain networks from behavioral labels.

\paragraph{What reasoning text can establish.}
Chain-of-thought explanations can omit influences on answers \citep{turpin2023faithfulness,chen2025faithfulness}. Relatedly, neuroimaging requires care when inferring cognitive processes from observed activity \citep{poldrack2006inference}. Our signatures characterize emitted text: decoding task identity does not identify its mechanism, and changing annotations alongside accuracy does not establish mediation. Multiple outcome protocols constrain interpretation.
